% ============================================================
%  Báo cáo Kỹ thuật: Mô phỏng NR Cell 128T128R Massive MIMO
%  với Codebook 3GPP Release 19 eTypeII
%
%  Tác giả : ThangTQ23 — VSI, 04/2026
%  Compile : pdflatex × 2  (hoặc upload lên Overleaf)
% ============================================================
\documentclass[12pt, a4paper]{article}

% ── Packages ─────────────────────────────────────────────────
\usepackage[T5]{fontenc}
\usepackage[utf8]{inputenc}
\usepackage[vietnamese,english]{babel}
\usepackage{lmodern}
\usepackage{geometry}
\geometry{left=2.5cm, right=2.5cm, top=2.5cm, bottom=2.5cm}

\usepackage{amsmath, amssymb}
\usepackage{booktabs}
\usepackage{longtable}
\usepackage{array}
\usepackage{multirow}
\usepackage{graphicx}
\usepackage{xcolor}
\usepackage{listings}
\usepackage{hyperref}
\usepackage{enumitem}
\usepackage{caption}
\usepackage{fancyhdr}
\usepackage{titlesec}
\usepackage{microtype}
\usepackage{tabularx}
\usepackage{ragged2e}

% ── Hyperref ─────────────────────────────────────────────────
\hypersetup{
  colorlinks = true,
  linkcolor  = blue!60!black,
  urlcolor   = blue!70!black,
  pdftitle   = {Bao cao Ky thuat: NR Cell 128T128R (Rel-19)},
  pdfauthor  = {ThangTQ23 -- VSI}
}

% ── Code listing style ───────────────────────────────────────
\lstdefinestyle{matlab}{
  language        = Matlab,
  basicstyle      = \ttfamily\footnotesize,
  keywordstyle    = \color{blue}\bfseries,
  commentstyle    = \color{green!50!black}\itshape,
  stringstyle     = \color{red!70!black},
  numbers         = left,
  numberstyle     = \tiny\color{gray},
  stepnumber      = 1,
  numbersep       = 6pt,
  backgroundcolor = \color{gray!8},
  frame           = single,
  frameround      = tttt,
  breaklines      = true,
  showstringspaces= false,
  tabsize         = 4
}
\lstset{style=matlab}

% ── Header / Footer ──────────────────────────────────────────
\pagestyle{fancy}
\fancyhf{}
\lhead{\small 128T128R NR Cell --- Release 19 eTypeII}
\rhead{\small Báo cáo Nội bộ, 04/2026}
\cfoot{\thepage}
\renewcommand{\headrulewidth}{0.4pt}

% ── Section formatting ───────────────────────────────────────
\titleformat{\section}{\large\bfseries\color{blue!60!black}}{}{0em}{\thesection\quad}
\titleformat{\subsection}{\normalsize\bfseries}{}{0em}{\thesubsection\quad}
\titleformat{\subsubsection}{\normalsize\itshape\bfseries}{}{0em}{\thesubsubsection\quad}

% ── Custom commands ──────────────────────────────────────────
\newcommand{\spec}[1]{\textcolor{green!40!black}{\texttt{[#1]}}}
\newcommand{\file}[1]{\texttt{#1}}
\newcommand{\param}[1]{\texttt{#1}}
% Cột bảng căn trái, wrap text
\newcolumntype{L}[1]{>{\RaggedRight\arraybackslash}p{#1}}

% ============================================================
\begin{document}

% ── Trang bìa ────────────────────────────────────────────────
\begin{titlepage}
  \centering
  \vspace*{2cm}
  {\LARGE\bfseries Báo cáo Kỹ thuật\\[0.5em]
   Mô phỏng NR Cell 128T128R\\[0.3em]
   với Codebook 3GPP Release~19 CSI-RS
  \par}
  \vspace{1.5cm}
  {\large Báo cáo Nội bộ \par}
  \vspace{0.5cm}
  {\large Tháng 4, 2026 \par}
  \vspace{1.5cm}
  \rule{0.8\textwidth}{0.4pt}\\[0.6em]
  \begin{tabular}{ll}
    \textbf{Tác giả}    & ThangTQ23, VSI \\
    \textbf{Phiên bản}  & v1.0 \\
    \textbf{Trạng thái} & Sẵn sàng bàn giao cho DRL training \\
    \textbf{Ngày}       & 16/04/2026 \\
  \end{tabular}\\[0.6em]
  \rule{0.8\textwidth}{0.4pt}
  \vspace{1.8cm}
  \begin{abstract}
    \noindent
    Báo cáo này ghi lại quá trình thiết kế, hiện thực và kiểm chứng môi
    trường mô phỏng downlink 128T128R massive MIMO tuân thủ 3GPP Release~19 sử dụng CSI-RS,
    xây dựng trên nền MATLAB Wireless Network Toolbox và 5G Toolbox.
    Xuất phát từ baseline link-level Rel-15/16, dự án đã trải qua 4 phiên
    bản để cho ra môi trường NR~Cell system-level đóng vai trò gym cho DRL
    agent học policy scheduling MU-MIMO.
    Đóng góp chính là việc tích hợp các refined codebook (Type I Mode A/B và eTypeII) Release~19 (TS~38.214 \S5.2.2.2.5a) hỗ trợ 128~port, cho phép mô phỏng CSI
    feedback chính xác về mặt không gian cho massive MIMO.
    Báo cáo cung cấp bằng chứng kiểm chứng, bảng tra cứu đặc tả kỹ thuật
    và hướng dẫn sử dụng đầy đủ phục vụ bàn giao dự án.
  \end{abstract}
\end{titlepage}

\tableofcontents
\newpage

% ============================================================
\section{Giới thiệu và Động lực}
% ============================================================

Dự án được khởi tạo để giải quyết khoảng trống trong bộ công cụ mô
phỏng 5G hiện tại: MATLAB Wireless Network Toolbox tiêu chuẩn chỉ hỗ
trợ tối đa 32~anten phát.

Công việc trong báo cáo này giải quyết:

\begin{enumerate}[nosep]
  \item \textbf{Benchmark link-level} (v1/v2): so sánh bốn phương pháp codebook
        từ Rel-15/16 Type~I đến Rel-19 eTypeII trên kênh 128~port massive MIMO.
  \item \textbf{Bộ mô phỏng NR Cell system-level} (v3): ghép cặp
        MU-MIMO, thích nghi liên kết OLLA.
  \item \textbf{Môi trường gym cho DRL} (v4): giao diện TCP socket, observation spaces với 8 features và một chương trình baseline để so sánh các thuật toán scheduler (RR, PF, BestCQI).
\end{enumerate}

Việc nâng cấp lên Release~19 là đóng góp cốt lõi: cho phép mô phỏng
CSI feedback beamforming tiên tiến nhất theo chuẩn NR~Rel-19
(3GPP TR~38.214), tạo môi trường huấn luyện thực tế cho DRL scheduler
hướng đến các triển khai thế hệ tiếp theo.

% ============================================================
\section{Tiến trình Phát triển và Cấu trúc File}
% ============================================================

\subsection{Lịch sử phiên bản}

\begin{center}
\begin{tabular}{L{1.4cm} L{6.1cm} L{7cm}}
\toprule
\textbf{Phiên bản} & \textbf{File} & \textbf{Mô tả} \\
\midrule
v1.0 & \file{CSIRS\_128T128R\_v1\_0.m} &
  Link-level, 4 approaches (A/B/C/D), 1~UE, 1 channel realization \\[3pt]
v2.0 & \file{CSIRS\_128T128R\_v2\_0.m} &
  SNR sweep + Monte Carlo (B/C/D/E), subband PMI/CQI \\[3pt]
v3.0 & \file{CSIRS\_128T128R\_v3\_0\_MUMIMO.m} &
  System-level NR Cell, mô phỏng MU-MIMO event-driven \\[3pt]
v4.0-DRL & \file{CSIRS\_128T128R\_v4\_0\_SOCKET.m} &
  DRL gym: giao thức TCP socket, FTP-3, observation spaces: 8 features \\[3pt]
v4.0-BL & \file{CSIRS\_128T128R\_v4\_0\_BASELINE.m} &
  Môi trường giống v4.0-DRL, scheduler baseline so sánh \\
\bottomrule
\end{tabular}
\captionof{table}{Lịch sử phiên bản và các file mô phỏng chính.}
\label{tab:versions}
\end{center}

\subsection{Các module hỗ trợ}

\begin{center}
\begin{tabular}{L{5cm} L{9.5cm}}
\toprule
\textbf{File} & \textbf{Vai trò} \\
\midrule
\file{csirs\_feedback.m} &
  Engine CSI feedback: RI/PMI/CQI cho Approaches B, C, D, E \\[3pt]
\file{csirs\_computeMetrics.m} &
  Ánh xạ SINR $\to$ MCS $\to$ throughput \\[3pt]
\file{csirs\_buildGrid.m} &
  Xây dựng TX resource grid (4 tài nguyên NZP-CSI-RS) \\[3pt]
\file{csirs\_channelEstimate.m} &
  Ước lượng kênh MMSE trên 2~slot \\[3pt]
\file{csirs\_channelSetup.m} &
  Khởi tạo đối tượng kênh CDL \\[3pt]
\file{csirs\_txRx.m} &
  Chuỗi TX/kênh/AWGN/giải điều chế \\[3pt]
\file{csirs\_printSweepTable.m} &
  Định dạng kết quả bảng SNR sweep \\[3pt]
\file{nrDRLScheduler.m} &
  \texttt{classdef} kế thừa \texttt{nrScheduler}: inject DRL action + socket \\[3pt]
\file{test\_dummy\_server.py} &
  Server Python kiểm thử (Level~1--3 validation) \\[3pt]
\file{hNRCreateCDLChannels.m} &
  Tạo kênh CDL cho nhiều UE \\[3pt]
\file{hNRCustomChannelModel.m} &
  Adapter kênh tùy chỉnh cho \texttt{wirelessNetworkSimulator} \\
\bottomrule
\end{tabular}
\captionof{table}{Các module hỗ trợ và vai trò.}
\label{tab:modules}
\end{center}

% ============================================================
\section{Nâng cấp 3GPP Release~19: Chi tiết Kỹ thuật}
% ============================================================

\subsection{Tổng quan thay đổi so với Baseline Rel-15/16}

Approach~A ban đầu sử dụng codebook Rel-15/16 Type~I Single-Panel
cấu hình cho panel \emph{ảo} 32~port (TS~38.214 \S5.2.2.2.1). Nâng
cấp Release~19 thay thế bằng ba chế độ codebook mới vận hành trên
\emph{toàn bộ} mảng 128~port:

\begin{enumerate}
  \item \textbf{Approach C} — Rel-19 Type~I Single-Panel Mode~A
        \spec{TS~38.214 \S5.2.2.2.1a, CodebookMode=1}
  \item \textbf{Approach D} — Rel-19 Type~I Single-Panel Mode~B
        \spec{TS~38.214 \S5.2.2.2.1a, CodebookMode=2}
  \item \textbf{Approach E} — Rel-19 Refined eTypeII 128~port
        \spec{TS~38.214 \S5.2.2.2.5a}
\end{enumerate}

\subsection{Cấu hình Anten}

Mảng anten 128~port được cấu hình dạng UPA (Uniform Planar Array)
phân cực kép 2D:
\[
  \text{Mảng gNB} = 4V \times 16H \times 2\,\text{pol} = 128\text{ port}
\]
Theo ký hiệu 3GPP: $N_g = 1$ panel, $N_1 = 16$ (chiều ngang),
$N_2 = 4$ (chiều đứng), kích thước panel $[N_g, N_1, N_2] = [1, 16, 4]$.

Tham số MATLAB tương ứng:
\begin{lstlisting}
panelDimensions = [1, 16, 4];       % [Ng=1, N1=16, N2=4] -- TS 38.214 S5.2.2.2.5a
gnbArraySize    = [4, 16, 2, 1, 1]; % [Nv, Nh, Npol, Mg, Mg2] = 128 port
\end{lstlisting}

\subsection{Cấu hình CSI-RS}

Bốn tài nguyên NZP-CSI-RS (Row~18, CDM8) được dùng để đo toàn bộ
128~port trên 2~slot liên tiếp:
\begin{itemize}[nosep]
  \item 4 tài nguyên $\times$ 32~port mỗi tài nguyên = 128~port tổng
  \item Row~18 hỗ trợ 32~port, CDM length~8 \spec{TS~38.211 \S7.4.1.5}
  \item Các tài nguyên phân bổ trên slot~0 và slot~1 để tránh chồng chéo
\end{itemize}

\subsection{Rel-19 Type~I Refined Single-Panel (Approaches C và D)}

Cả Approach~C lẫn Approach~D đều dùng
\param{CodebookType = 'typeI-SinglePanel-r19'}, kích hoạt đường xử lý
Rel-19 trong \texttt{nr5g.internal.nrRISelect} và
\texttt{nr5g.internal.nrCQISelect} của MATLAB.

\begin{itemize}
  \item \textbf{Mode~A} (\texttt{CodebookMode=1}): hệ số oversampling
        $O_1 = O_2 = 4$, beam+co-phase theo subband. PMI và CQI đều
        theo \param{pmiMode}.
  \item \textbf{Mode~B} (\texttt{CodebookMode=2}): nhóm beam ($i_1$)
        luôn là wideband (ràng buộc đặc tả: tìm kiếm nhóm beam trên
        toàn băng thông); co-phase ($i_2$) có thể theo subband.
\end{itemize}

\subsection{Rel-19 Refined eTypeII --- Approach E}
\label{sec:etypeii}

Đây là đóng góp cốt lõi của nâng cấp Release~19. Codebook eTypeII-r19
(TS~38.214 \S5.2.2.2.5a) chọn $L$ beam DFT cơ sở mỗi phân cực và kết
hợp chúng với hệ số biên độ+pha phức:
\[
  \mathbf{W} = \mathbf{W}_1 \mathbf{W}_2, \quad
  \mathbf{W}_1 \in \mathbb{C}^{N_T \times 2L}, \quad
  \mathbf{W}_2 \in \mathbb{C}^{2L \times r}
\]
trong đó $\mathbf{W}_1$ là ma trận chọn beam DFT khối chéo (beam dùng
chung cho cả hai phân cực) và $\mathbf{W}_2$ chứa hệ số tổ hợp tuyến
tính được lượng tử hóa.

\subsubsection{Lựa chọn ParameterCombination}

Trường \param{ParameterCombination} (PC) kiểm soát số beam $L$, độ
phân giải biên độ $p_v$ và overhead phân số $\beta$:

\begin{center}
\begin{tabular}{ccccL{5cm}}
\toprule
\textbf{PC} & $L$ & $p_v$ & $\beta$ & \textbf{Ghi chú} \\
\midrule
1 & 2 & 1/2 & 1/4 & Overhead thấp, rank 1--2 \\
2 & 2 & 1   & 1/2 & Biên độ đầy đủ, rank 1--2 \\
3 & 4 & 1/4 & 1/4 & Overhead trung bình \\
4 & 4 & 1/2 & 1/4 & Overhead chuẩn \\
\textbf{5} & \textbf{4} & \textbf{1/4} & \textbf{1/4} &
  \textbf{Dùng trong dự án này} \\
6 & 4 & 1/2 & 1/2 & Độ phân giải biên độ cao hơn \\
7 & 6 & ---  & --- & Rank tối đa = 2 \\
8 & 6 & ---  & --- & Rank tối đa = 2 \\
\bottomrule
\end{tabular}
\captionof{table}{Các giá trị ParameterCombination theo TS~38.214
Table~5.2.2.2.5-1. PC=5 ($L=4$, $\beta=1/4$) được dùng trong mô
phỏng system-level.}
\label{tab:pc}
\end{center}

\subsubsection{Hai đường hiện thực}

Do bùng nổ tổ hợp khi $L \geq 4$, hai đường được sử dụng:
\begin{itemize}
  \item \textbf{$L = 2$ (PC 1/2):} Dùng toàn bộ pipeline toolbox qua
        \texttt{nr5g.internal.nrPMIReport} tìm kiếm toàn diện trên
        $\binom{1024}{2} \approx 31.000$ cặp beam --- khả thi về bộ
        nhớ, cho $i_1 = [i_{1,1},\, i_{1,2},\, i_{1,\mathrm{LC}}]$
        đúng theo đặc tả.
  \item \textbf{$L \geq 4$ (PC 3--8):} Greedy DFT matching pursuit
        (hiện thực trong \file{csirs\_feedback.m}:
        \texttt{greedyEType2R19Precoder}) --- tránh OOM từ
        $\binom{1024}{4} \times 128 \times 4 > 5\,\text{GB}$, giữ
        nguyên cấu trúc $\mathbf{W}_1\mathbf{W}_2$ theo đặc tả.
\end{itemize}

\subsection{Override Thư viện MATLAB --- Kích hoạt 128~Port và Rel-19}

MATLAB 5G Toolbox tiêu chuẩn hỗ trợ tối đa 32~anten phát và không
cung cấp các kiểu codebook Rel-19 (\param{'eTypeII-r19'},
\param{'typeI-SinglePanel-r19'}) qua API công khai.

Để bổ sung các khả năng này, các phiên bản \emph{đã sửa} của
thư viện \texttt{5g/} và \texttt{wirelessnetwork/} được thêm vào
MATLAB path:
\begin{lstlisting}
addpath(genpath(fullfile(projectRoot, 'wirelessnetwork')), '-begin');
addpath(genpath(fullfile(projectRoot, '5g')),              '-begin');
\end{lstlisting}

\begin{enumerate}[nosep]
  \item \texttt{nrGNB} với \param{NumTransmitAntennas = 128}
        (trước đây giới hạn 32)
  \item Đường xử lý codebook eTypeII-r19 trong \texttt{nrUEAbstractPHY}
  \item Hỗ trợ \param{'typeI-SinglePanel-r19'} và \param{'eTypeII-r19'}
        trong \texttt{nr5g.internal.nrRISelect} / \texttt{nrCQISelect}
  \item \texttt{nrNodeValidation} chấp nhận \param{PMICQIMode = "Subband"}
        cho cấu hình gNB 128~port
\end{enumerate}

% ============================================================
\section{Bảng Tra cứu Đặc tả Kỹ thuật}
% ============================================================

Bảng~\ref{tab:spec} ánh xạ từng tham số cấu hình chính đến tham chiếu
specs 3GPP tương ứng, cung cấp audit trail trực tiếp từ code đến
tiêu chuẩn.

\begin{center}
\begin{longtable}{L{4.5cm} L{3.8cm} L{5.5cm}}
\toprule
\textbf{Tham số / Tính năng} & \textbf{Giá trị (dự án này)} &
\textbf{Tham chiếu 3GPP} \\
\midrule
\endhead
\bottomrule
\endfoot
Kiểu codebook (Approach E)  & \texttt{'eTypeII-r19'}         &
  TS~38.214 \S5.2.2.2.5a \\[2pt]
Kích thước panel             & $[N_g,N_1,N_2]=[1,16,4]$       &
  TS~38.214 \S5.2.2.2.5a, Table~5.2.2.2.5-1 \\[2pt]
ParameterCombination         & PC=5 ($L=4$, $\beta=1/4$)      &
  TS~38.214 Table~5.2.2.2.5-1, dòng 5 \\[2pt]
Kiểu codebook (Approach C)   & \texttt{'typeI-SP-r19'}, Mode~A &
  TS~38.214 \S5.2.2.2.1a \\[2pt]
Kiểu codebook (Approach D)   & \texttt{'typeI-SP-r19'}, Mode~B &
  TS~38.214 \S5.2.2.2.1a \\[2pt]
Kiểu codebook (Approach A)   & \texttt{'typeI-SP'}, Rel-15/16  &
  TS~38.214 \S5.2.2.2.1 \\[2pt]
Oversampling DFT             & $O_1 = O_2 = 4$                &
  TS~38.214 \S5.2.2.2.3 \\[2pt]
CSI-RS row (32~port)         & Row~18, CDM8                   &
  TS~38.211 Table~7.4.1.5.3-1 \\[2pt]
Kích thước subband           & 4~RB (NRB $\in [24,72]$)       &
  TS~38.214 Table~5.2.1.4-2 \\[2pt]
Bảng CQI                     & Table~1                        &
  TS~38.214 Table~5.2.2.1-2 \\[2pt]
OLLA target BLER             & $\approx 10\%$                 &
  StepDown/(StepDown+StepUp) = 0.03/0.30 \\[2pt]
Chu kỳ SRS (UE)              & 20~ms                          &
  TS~38.211 \S6.4.1.4 \\[2pt]
Chu kỳ báo cáo CSI           & 10~slot ($= 5$~ms @ 30~kHz)   &
  TS~38.214 \S5.2.1 \\[2pt]
Tần số sóng mang             & 4,9~GHz (băng n79)             &
  TS~38.104 Table~5.4.2.3-1 \\[2pt]
NRB (10~MHz, SCS 30~kHz)     & 24                             &
  TS~38.104 Table~5.3.2-1 \\[2pt]
Hồ sơ trễ CDL                & CDL-D (LOS-dominant)           &
  TS~38.901 Table~7.7.2-4 \\[2pt]
Chế độ ghép song công        & TDD                            &
  TS~38.104 \S4.3 \\
\end{longtable}
\captionof{table}{Bảng tra cứu tham số — specs 3GPP.}
\label{tab:spec}
\end{center}

% ============================================================
\section{Bằng chứng Kiểm chứng}
% ============================================================

\subsection{Bằng chứng 1 --- Cấu trúc Chỉ số PMI đặc trưng cho Codebook}

Bằng chứng trực tiếp nhất cho thấy đường xử lý codebook eTypeII-r19
Rel-19 được kích hoạt đúng là cấu trúc của chỉ số PMI $i_1$:

\begin{itemize}
  \item \textbf{Rel-15/16 eTypeII:} $i_1 = [i_{1,1},\, i_{1,2}]$
        --- 2 phần tử (chỉ số cặp beam)
  \item \textbf{Rel-19 eTypeII-r19:} $i_1 = [i_{1,1},\, i_{1,2},\, i_{1,\mathrm{LC}}]$
        --- 3 phần tử, trong đó $i_{1,\mathrm{LC}}$ là chỉ số hệ số
        tổ hợp tuyến tính \spec{TS~38.214 \S5.2.2.2.5a}
\end{itemize}

Test Level~0 (chạy với \file{CSIRS\_128T128R\_v3\_0\_MUMIMO.m}) đã
xác nhận:
\[
  i_1 = [4,\; 2,\; 1431]
\]
Cấu trúc 3~phần tử này là điều \emph{không thể} đạt được từ bất kỳ
cấu hình codebook Rel-15/16 nào, xác định một cách rõ ràng rằng
đường code Rel-19 đang được sử dụng.

\subsection{Bằng chứng 2 --- Bất đẳng thức Phân cấp Hiệu năng}

Bốn approach tạo thành một phân cấp hiệu năng có cơ sở lý thuyết
\emph{không thể vi phạm} nếu mỗi codebook được hiện thực đúng:
\[
  \underbrace{\text{B (SVD upper bound)}}_{\text{lý tưởng}} \;\geq\;
  \underbrace{\text{E (eTypeII-r19)}}_{L=4\text{ beam}} \;\geq\;
  \underbrace{\text{C (TypeI-r19 ModeA)}}_{\text{lưới codebook}} \;\geq\;
  \underbrace{\text{A (TypeI Rel-15/16)}}_{\text{chỉ 32 port}}
\]

\textbf{Tại sao đây là bằng chứng:}
\begin{itemize}
  \item B là SVD upper bound --- mọi precoder dựa trên feedback
        \emph{phải} thực hiện tại hoặc dưới ngưỡng này theo cấu trúc.
  \item E dùng $L=4$ beam DFT với cấu trúc $\mathbf{W}_1\mathbf{W}_2$
        tổng quát hóa chặt chẽ codebook TypeI (TypeI là trường hợp đặc
        biệt $L=1$ beam, không có tổ hợp tuyến tính).
  \item A hoạt động trên các mảng con 32~port ảo, bỏ qua thông tin
        không gian liên mảng con có trong kênh 128~port đầy đủ.
\end{itemize}

Nếu phân cấp được thỏa mãn trên toàn bộ điểm SNR và các lần thực
hiện kênh trong quét v2 (\file{CSIRS\_128T128R\_v2\_0.m}), tính đúng
đắn được xác nhận. Mức tăng tương đối kỳ vọng tại SNR cao (25~dB)
theo 3GPP TR~38.859:
\begin{itemize}[nosep]
  \item E so với C: $+15\text{--}25\%$ thông lượng
  \item Khoảng cách E so với B: $5\text{--}15\%$ (overhead lượng tử hóa)
\end{itemize}

\subsection{Bằng chứng 3 --- Hội tụ OLLA trong Mô phỏng System-level}

Outer Loop Link Adaptation (OLLA) hội tụ BLER từng UE về giá trị
mục tiêu $\approx 10\%$ theo tiến trình mô phỏng. Mục tiêu được suy
ra từ kích thước bước:
\[
  \mathrm{BLER}_{\mathrm{target}} = \frac{\Delta_{\mathrm{down}}}{\Delta_{\mathrm{down}} + \Delta_{\mathrm{up}}} = \frac{0{,}03}{0{,}03 + 0{,}27} = 10\%
\]
Sự hội tụ OLLA (quan sát từ frame~20 trở đi trong các lần chạy 50~frame)
xác nhận: (a)~CSI feedback đang được nhận và xử lý, (b)~tín hiệu HARQ
ACK/NACK hoạt động đúng, và (c)~lựa chọn MCS đang phản hồi theo
feedback. Đây là kiểm thử tích hợp end-to-end.

\subsection{Bằng chứng 4 --- Kiểm thử Giao thức Socket DRL}

Bốn cấp kiểm thử đã được thực hiện để kiểm chứng môi trường DRL v4:

\begin{center}
\begin{tabular}{cL{4.5cm}L{3.7cm}c}
\toprule
\textbf{Level} & \textbf{Mô tả} & \textbf{Công cụ} & \textbf{Kết quả} \\
\midrule
0 & Chạy baseline v3 (5 frame), kiểm tra cấu trúc i1 &
  \file{v3\_0\_MUMIMO.m} & PASS \\[3pt]
1 & Handshake socket + giao thức TTI\_START/DONE &
  \file{test\_dummy\_server.py} (NOOP) & PASS \\[3pt]
2 & Kiểm tra giá trị quan sát (JSON dump, 8~đặc trưng) &
  \file{test\_dummy\_server.py} (obs\_dump) & PASS \\[3pt]
3 & Smoke test hành động ngẫu nhiên, thông lượng cell~$> 0$ &
  \file{test\_dummy\_server.py} (random) & PASS \\
\bottomrule
\end{tabular}
\captionof{table}{Các cấp kiểm thử môi trường DRL. Cả bốn cấp đều đạt.}
\label{tab:tests}
\end{center}

% ============================================================
\section{Kiến trúc System-level (v3/v4)}
% ============================================================

\subsection{Framework Mô phỏng}

Mô phỏng system-level sử dụng framework event-driven
\texttt{wirelessNetworkSimulator} của MATLAB (Wireless Network Toolbox).
Các thuộc tính chính:

\begin{itemize}[nosep]
  \item \textbf{Nút:} 1 gNB (128T128R) + 16~UE (4Tx/4Rx mỗi UE)
  \item \textbf{Lưu lượng:} FTP Model~3 (6~Mbps, gói 1500~byte) cho
        động học buffer; Full Buffer tùy chọn cho phân tích trạng thái ổn định
  \item \textbf{Kênh:} CDL-D, RMS delay spread 100~ns,
        Doppler 5~Hz ($\approx$1~km/h @ 4,9~GHz)
  \item \textbf{Scheduler:} RAT-0 (cấp phát tài nguyên theo RBG)
\end{itemize}

\subsection{Cấu hình MU-MIMO}

\begin{center}
\begin{tabular}{L{5.5cm} L{8cm}}
\toprule
\textbf{Tham số} & \textbf{Giá trị / Giải thích} \\
\midrule
MaxNumUsersPaired       & 4 \\[3pt]
MaxNumLayers            & 8 \\[3pt]
SemiOrthogonalityFactor & 0,7 \\[3pt]
MinCQI                  & 1 \\[3pt]
Tín hiệu đo CSI         & CSI-RS \\[3pt]
Cơ sở ghép cặp          & $L=2$ beam $\times$ 2~pol = không gian con 4D;
                          $K \cdot r \leq 4$ để đảm bảo tính trực giao \\
\bottomrule
\end{tabular}
\captionof{table}{Cấu hình scheduler MU-MIMO.}
\end{center}

\subsection{Triển khai UE}

UE được triển khai trong sector 120° với:
\begin{itemize}[nosep]
  \item Góc phương vị: đều trong $[-60°, +60°]$
  \item Khoảng cách: $r \in [50, 500]$~m, lấy mẫu từ phân phối đều
        trên \emph{diện tích} sector qua ``sqrt-trick'':
        $r = \sqrt{r_{\min}^2 + u(r_{\max}^2 - r_{\min}^2)}$,
        $u \sim \mathcal{U}[0,1]$
\end{itemize}

\subsection{Cấu hình OLLA}

OLLA (State~3: không reset) tích lũy MCS offset xuyên suốt mô phỏng:
\begin{lstlisting}
ollaConfig = struct('InitialOffset', 0, 'StepUp', 0.27, 'StepDown', 0.03);
% Target BLER = 0.03 / (0.03 + 0.27) = 10%
% State 3: reset da duoc xoa -- thiet ke co chu dich
\end{lstlisting}

Thiết kế không reset là chủ đích: CQI MU-MIMO được đo theo SU nên
đánh giá quá cao SINR. OLLA phải tích lũy qua nhiều grant để học đúng
mức giảm MCS trong điều kiện MU-MIMO.

% ============================================================
\section{Môi trường Gym DRL (v4)}
% ============================================================

\subsection{Kiến trúc tổng thể}

\begin{verbatim}
Python DRL Agent  <-- TCP socket :5555 -->  MATLAB NR Cell (v4_0_SOCKET)
  PPO / SAC / ...                            wirelessNetworkSimulator
       |                                            |
  action: [NumLayers x NumRBGs]            obs: TTI_START JSON
  reward: PF throughput                    nrDRLScheduler (classdef)
                                           eTypeII-r19, subband CQI
\end{verbatim}

\subsection{Observation spaces: 8 features}

\begin{center}
\begin{tabular}{cL{2.7cm}L{5.5cm}L{3.5cm}}
\toprule
\textbf{\#} & \textbf{JSON Key} & \textbf{Mô tả} & \textbf{Kích thước} \\
\midrule
1 & \texttt{avg\_tp}       & EMA thông lượng DL (Mbps)              & $[\text{MaxUEs}]$ \\[3pt]
2 & \texttt{ue\_rank}      & Rank indicator mỗi UE                  & $[\text{MaxUEs}]$ \\[3pt]
3 & \texttt{occupied\_rbgs}& Bản đồ cấp phát RBG                    & theo dõi bởi Python \\[3pt]
4 & \texttt{buf}           & DL buffer (bytes)                      & $[\text{MaxUEs}]$ \\[3pt]
5 & \texttt{curr\_mcs}     & MCS hiệu dụng (sau OLLA)               & $[\text{MaxUEs}]$ \\[3pt]
6 & \texttt{ue\_i1}        & Chỉ số beam-group $i_1$ eTypeII-r19    & $[\text{MaxUEs} \times 3]$ \\[3pt]
7 & \texttt{sub\_cqi}      & Subband CQI (resample sang lưới RBG)   & $[\text{MaxUEs} \times \text{nRBG}]$ \\[3pt]
8 & \texttt{max\_cross\_corr} & $\kappa$ tương quan chéo max-cột (per-RBG) &
  $[\text{MaxUEs}^2 \times \text{nRBG}]$ \\
\bottomrule
\end{tabular}
\captionof{table}{Các đặc trưng không gian quan sát gửi đến DRL agent mỗi TTI.}
\label{tab:obs}
\end{center}

\noindent\textbf{Lưu ý hiện thực quan trọng:}
\begin{itemize}[nosep]
  \item \texttt{avg\_tp} dùng EMA (\texttt{ActualTputEMA}), không phải
        thông lượng tức thời
  \item \texttt{eligible\_ues} gửi 0-based (MATLAB RNTI~$-$~1) để khớp
        với index mảng Python
  \item \texttt{max\_cross\_corr} được định nghĩa là
        $\kappa(W_i, W_j) = \max_k |W_i[:,k]^\dagger W_j[:,k]| / (\|w_i^k\|\|w_j^k\|)$,
        tính per-RBG --- \emph{không phải} hàm semi-orthogonality nội bộ của MATLAB
\end{itemize}

\subsection{Action spaces}

DRL agent trả về ma trận $A \in \mathbb{Z}^{N_L \times M}$ trong đó:
\begin{itemize}[nosep]
  \item $N_L$ = số layer không gian (MU-MIMO stream)
  \item $M$ = số RBG ($= \lceil N_{\rm RB} / \text{kích thước RBG} \rceil$)
  \item $A[l,m] \in \{1, \ldots, |\mathcal{U}|\}$: chỉ số UE (1-based) cho layer~$l$, RBG~$m$
  \item $A[l,m] = |\mathcal{U}|+1$: không cấp phát (skip)
\end{itemize}

\subsection{Giao thức Socket}

\begin{verbatim}
HANDSHAKE  ->  "HELLO\n"
           <-  "READY\n"

Mỗi TTI:
  ->  "TTI_START " + JSON(obs) + "\n"
  <-  "ACTION " + JSON(action) + "\n"
  ->  "TTI_DONE " + JSON(reward) + "\n"

Kết thúc episode:
  ->  "STOP\n"
  <-  (kết nối đóng)
\end{verbatim}

% ============================================================
\section{Hướng dẫn Sử dụng}
% ============================================================

\subsection{Yêu cầu Hệ thống}

\begin{itemize}[nosep]
  \item MATLAB R2024a trở lên (yêu cầu 5G Toolbox + Wireless Network Toolbox)
  \item Thư mục thư viện đã sửa \file{5g/} và \file{wirelessnetwork/}
        có mặt tại thư mục gốc dự án (đã bao gồm trong repository)
  \item Python 3.9+ với \texttt{socket} (thư viện chuẩn) cho giao diện DRL
\end{itemize}

\subsection{So sánh Codebook Link-level (v2)}

\begin{lstlisting}
% Chay tu MATLAB command window, CWD la thu muc goc du an
% Cho ra bang sweep SNR: RI / MCS / Throughput cho B, C, D, E

run('CSIRS_128T128R_v2_0.m')

% Tham so co the chinh (dau file):
%   pmiMode    = 'Subband' | 'Wideband'
%   cqiMode    = 'Wideband' | 'Subband'
%   paramCombE = 5       (PC cho eTypeII, xem Bang 3)
%   snrList    = [-5, 10, 25]  (dB)
%   nRealiz    = 20      (so lan thuc hien kenh moi diem SNR)
\end{lstlisting}

\textbf{Kết quả kỳ vọng:} Bảng in ra terminal hiển thị RI, MCS và
thông lượng cho B/C/D/E tại từng điểm SNR. Phân cấp
B $\geq$ E $\geq$ C/D phải thỏa mãn tại mọi điểm.

\subsection{Mô phỏng NR Cell System-level (v3)}

\begin{lstlisting}
% Chay tu MATLAB command window
% Yeu cau Wireless Network Toolbox

run('CSIRS_128T128R_v3_0_MUMIMO.m')

% §8:  in tom tat dieu kien mo phong
% §9:  chay 5 frame (numFrameSimulation = 5)
% §10: KPI Summary (displayPerformanceIndicators)
% §11: luu simulationLogs.mat
% §11b: bang Grant/MCS tung UE
% §11c: xu huong BLER & MCS tung frame (OLLA Convergence Monitor)
% §12: histogram ghep cap MU-MIMO

% Thay doi che do PMI/CQI (dau file):
%   pmiCQIMode = "Subband";   % khuyen nghi
%   pmiCQIMode = "Wideband";  % nhanh hon khi debug
\end{lstlisting}

\subsection{Baseline Scheduler (v4 Baseline)}

\begin{lstlisting}
% Chay doc lap -- khong can Python server
% Cho ra KPI 50 frame lam baseline so sanh DRL

run('CSIRS_128T128R_v4_0_BASELINE.m')

% Tuy chon scheduler (dau file):
%   schedulerType = 'ProportionalFair'  % baseline chuan trong DRL
%   schedulerType = 'RoundRobin'        % can duoi
%   schedulerType = 'BestCQI'           % toi da hoa thong luong
\end{lstlisting}

\subsection{Môi trường DRL Socket (v4 Socket)}

\textbf{Bước 1:} Khởi động server Python (chế độ kiểm thử hoặc DRL agent thực):
\begin{lstlisting}[language=Python]
# Server kiem thu (Level 1: NOOP / Level 2: obs dump / Level 3: random)
python test_dummy_server.py --mode random --port 6666
\end{lstlisting}

\textbf{Bước 2:} Chạy mô phỏng MATLAB:
\begin{lstlisting}
run('CSIRS_128T128R_v4_0_SOCKET.m')
% Ket noi den localhost:6666 (DRL_Port trong nrDRLScheduler)
% Log ghi vao logs_v4/ kem timestamp
\end{lstlisting}

\textbf{Bước 3:} Với DRL agent thực, hiện thực giao thức socket mô
tả trong Mục~7.3 ở phía Python.

\subsection{Lưu ý Cấu hình Quan trọng}

\begin{enumerate}
  \item \param{gNB.PMICQIMode = pmiCQIMode} \textbf{phải} được gọi
        trước \texttt{connectUE(gNB, UEs, ...)}
  \item Dùng tên string cho scheduler: \param{'ProportionalFair'},
        \param{'RoundRobin'}, \param{'BestCQI'} (không dùng mã số)
  \item Không xóa các lệnh \texttt{addpath(..., '-begin')} --- nếu
        thiếu, MATLAB sẽ dùng toolbox mặc định và 128~port + eTypeII-r19
        sẽ lặng lẽ quay về trạng thái không được hỗ trợ
  \item OLLA \emph{không} reset giữa các frame --- đây là thiết kế
        có chủ đích (tích lũy State~3)
\end{enumerate}

% ============================================================
\section{Tổng kết Đóng góp}
% ============================================================

\begin{center}
\begin{tabular}{L{3.5cm} L{10.5cm}}
\toprule
\textbf{Hạng mục} & \textbf{Nội dung công việc} \\
\midrule
gNB 128~port & Sửa thư viện MATLAB để bổ sung
  \texttt{NumTransmitAntennas=128} và stack \texttt{nrUEAbstractPHY}
  tương ứng \\[3pt]
Rel-19 TypeI & Triển khai Approaches C và D dùng
  \texttt{'typeI-SinglePanel-r19'} (TS~38.214 \S5.2.2.2.1a)
  với cả Mode~A và Mode~B \\[3pt]
Rel-19 eTypeII & Triển khai Approach E: đường $L \leq 2$ dùng tìm kiếm
  toàn diện qua toolbox; đường $L \geq 4$ dùng greedy DFT matching
  pursuit (TS~38.214 \S5.2.2.2.5a) \\[3pt]
Subband PMI/CQI & Precoder và CQI per-subband
  (TS~38.214 Table~5.2.1.4-2) \\[3pt]
MU-MIMO System-level & NR Cell simulation với OLLA, ghép cặp MU-MIMO,
  kênh CDL-D, \texttt{wirelessNetworkSimulator} event-driven \\[3pt]
Giao diện DRL & \texttt{nrDRLScheduler} classdef với giao thức TCP
  socket, observation spaces: 8 features (đã kiểm chứng), FTP-3 \\[3pt]
Baseline & \texttt{v4\_0\_BASELINE.m} với ProportionalFair/RoundRobin/BestCQI
  cho so sánh DRL \\
\bottomrule
\end{tabular}
\captionof{table}{Tổng kết đóng góp của dự án.}
\end{center}

% ============================================================
\section{Kết luận}
% ============================================================

Dự án đã hiện thực và kiểm chứng thành công môi trường mô phỏng NR
Cell 128T128R massive MIMO tuân thủ 3GPP Release~19. Việc nâng cấp
từ codebook Rel-15/16 Type~I lên Rel-19 eTypeII được hỗ trợ bởi
bốn dòng bằng chứng độc lập: (1)~nhận dạng đặc trưng cấu trúc
chỉ số PMI, (2)~kiểm chứng phân cấp hiệu năng vật lý, (3)~hội tụ
OLLA trong mô phỏng system-level, và (4)~kiểm thử giao thức socket
DRL end-to-end. Môi trường sẵn sàng để bàn giao cho nhóm DRL training.

% ============================================================
\appendix
\section{Tóm tắt 5 Approaches}
% ============================================================

\begin{center}
\begin{tabular}{cL{2.5cm}L{5.5cm}L{3.2cm}}
\toprule
\textbf{Ký hiệu} & \textbf{Tên} & \textbf{Tham số chính} & \textbf{Đặc tả} \\
\midrule
A & Rel-15/16 Type~I SP &
  \texttt{'typeI-SP'}, \texttt{PanelDimensions=[8,2]}, per-32-port &
  TS~38.214 \S5.2.2.2.1 \\[3pt]
B & SVD Upper Bound &
  Full 128-port SVD, tham chiếu lý tưởng &
  --- \\[3pt]
C & Rel-19 Type~I Mode~A &
  \texttt{'typeI-SP-r19'}, \texttt{Mode=1}, \texttt{[1,16,4]} &
  TS~38.214 \S5.2.2.2.1a \\[3pt]
D & Rel-19 Type~I Mode~B &
  \texttt{'typeI-SP-r19'}, \texttt{Mode=2}, \texttt{[1,16,4]} &
  TS~38.214 \S5.2.2.2.1a \\[3pt]
E & Rel-19 eTypeII 128-port &
  \texttt{'eTypeII-r19'}, PC=5, $L=4$, \texttt{[1,16,4]} &
  TS~38.214 \S5.2.2.2.5a \\
\bottomrule
\end{tabular}
\captionof{table}{Năm approach codebook được hiện thực trong dự án.}
\end{center}

\section{Hàm MATLAB Nội bộ Chính}

\begin{center}
\begin{tabular}{L{5.5cm} L{8cm}}
\toprule
\textbf{Hàm} & \textbf{Vai trò trong dự án} \\
\midrule
\texttt{nr5g.internal.nrRISelect}  &
  Chọn RI cho Approaches C, D, E (L=2) \\[3pt]
\texttt{nr5g.internal.nrCQISelect} &
  Chọn CQI cho Approaches C, D, E (L=2) \\[3pt]
\texttt{nr5g.internal.nrPMIReport} &
  Báo cáo PMI (gọi nội bộ bởi hai hàm trên) \\[3pt]
\texttt{nrCSIRS}                   &
  Tạo tín hiệu CSI-RS \\[3pt]
\texttt{nrChannelEstimate}         &
  Ước lượng kênh MMSE \\[3pt]
\texttt{nrGNB}                     &
  Node gNB (đã vá: hỗ trợ 128T) \\[3pt]
\texttt{nrUE}                      &
  Node UE \\[3pt]
\texttt{wirelessNetworkSimulator}  &
  Bộ mô phỏng event-driven \\[3pt]
\texttt{configureScheduler}        &
  Cấu hình scheduler (chỉ dùng tên string) \\
\bottomrule
\end{tabular}
\captionof{table}{Các hàm MATLAB nội bộ chính và vai trò.}
\end{center}

\end{document}
